On rule-derived family trees, a 2-layer relational message-passing model reached filtered MRR \rhval{main/path-mp-2-hop/transductive/mrr/mean:3} on held-out triples and \rhval{main/path-mp-2-hop/inductive/mrr/mean:3} on a fresh graph with new entities, with no measurable transfer loss, while a tuned TransE reached \rhval{main/transe-reimplemented/transductive/mrr/mean:3} on held-out triples and was at chance on new entities. The inductive failure of TransE is largely a consequence of the evaluation set-up: fitting only the new entity vectors on the observed new graph restores it to \rhval{main/transe+fold-in-reimplemented/inductive/mrr/mean:3}, but it remains far below the path model. The path model's advantage is not unconditional: it needs two layers and training with the query edge removed, and when only a quarter of the new graph is observed it is not distinguishable from fold-in TransE. These conclusions hold for a small synthetic world with one generative process and five seeds; testing them on benchmarks with unknown rules, with tuned RotatE/ComplEx baselines and with distribution shift between graphs is the natural next step.
